\documentclass{report}
\usepackage[utf8]{inputenc}
\usepackage[spanish, es-noshorthands]{babel}
\usepackage{amsmath}
\usepackage{amssymb}

\begin{document}

\chapter*{Introducción}

El problema de la satisfacibilidad booleana (SAT) ocupa un lugar central
en las ciencias de la computación. Fue el primer problema que se demostró
NP-completo en 1971, un resultado de Stephen Cook~\cite{cook1971}. Por
tanto, todo problema de la clase NP se
reduce a él, y decidirlo en tiempo polinomial resolvería si $P = NP$. A la vez, el SAT tiene aplicaciones en la criptografía, el diseño y la
verificación de circuitos, los problemas de satisfacción de restricciones
(CSP) y la programación lógica~\cite{gu1997sat}.
Por eso interesa entender qué lo hace difícil y delimitar los casos que
se resuelven en tiempo polinomial, como 2-SAT, Horn-SAT y
XOR-SAT~\cite{schaefer1978}.

En una línea de investigación de la Facultad de Matemática y Computación
de la Universidad de La Habana, el SAT se aborda desde la teoría de
lenguajes formales. A cada fórmula booleana se le asocian un lenguaje
regular y una gramática. Las cadenas del lenguaje regular representan
las interpretaciones de sus variables que satisfacen la fórmula,
asumiendo que no hay variables repetidas, y la gramática impone la
consistencia: que las ocurrencias de una misma variable tomen el mismo
valor. La satisfacibilidad de la fórmula
equivale entonces a que la intersección de ambos lenguajes no sea vacía.
Así, el SAT se traduce en un problema sobre lenguajes formales, y la
dificultad de decidirlo se traslada al vacío de esa intersección.

Esta línea cuenta con antecedentes según el formalismo de la gramática de
consistencia. En la satisfacibilidad booleana libre del contexto (SATCF),
la consistencia se impone con una gramática libre del contexto, y el SAT
se decide en tiempo polinomial para las fórmulas con orden de ocurrencias
libre del contexto~\cite{fernandezarias2007}. En la satisfacibilidad
booleana de concatenación de rango simple (SATSRC), la consistencia se
impone con una gramática de concatenación de rango simple, que admite
cruces entre ocurrencias, y el SAT se decide en tiempo polinomial
mientras la aridad necesaria esté acotada~\cite{aguileralopez2016}. Las
gramáticas de concatenación de rango (RCG) son el siguiente formalismo de
esta línea.

En una RCG, los no terminales, llamados predicados, reconocen tuplas
de subcadenas en lugar de una sola, y la aridad de un predicado es la
cantidad de subcadenas de la tupla. Una RCG es no lineal cuando una
misma variable aparece más de una vez en el lado derecho de una regla,
y así se puede exigir que distintas subcadenas coincidan.

Con las RCG lineales se reconocen los mismos lenguajes que con las sRCG,
de modo que con ellas se modela la consistencia de las mismas fórmulas.
Para modelar la consistencia de más fórmulas, y resolver así más
instancias del SAT, hace falta la no linealidad de las RCG~\cite{boullier1998proposal}.

Al pasar de las gramáticas libres del contexto y las de concatenación
de rango simple a las RCG no se conserva la tratabilidad del vacío de la
intersección. Bertsch y Nederhof demostraron que el vacío de la
intersección de una RCG con un lenguaje regular arbitrario es
indecidible, y que, cuando el lenguaje regular es finito, decidir si la
intersección es no vacía es NP-completo~\cite{bertsch2001rcg}. El
lenguaje regular de las interpretaciones de una fórmula es finito, así
que este enfoque queda en el caso decidible.

Con la no linealidad de las
RCG se expresa la consistencia de cualquier fórmula,
sin la restricción de orden de SATCF y SATSRC, pero el problema queda
fuera de las subclases con vacío de la intersección polinomial conocidas.
Aun así, ese vacío se decide al enumerar las cadenas del autómata y
reconocer cada una con la gramática de consistencia, pero en tiempo
exponencial y con los dos formalismos por separado. Falta, entonces, una
representación de esa intersección sobre la que estudiar su vacío en un
solo formalismo.

Teniendo en cuenta lo anterior, surge la siguiente pregunta:
¿es posible expresar la intersección de una gramática de concatenación de
rango con un lenguaje regular finito mediante una sola gramática de
concatenación de rango, sobre la que estudiar el vacío que decide el SAT?

El objeto de estudio de este trabajo es la resolución del SAT mediante
formalismos de la teoría de lenguajes; dentro de ella, el campo de acción
es el vacío de la intersección de una gramática de concatenación de rango
con un lenguaje regular finito.

El objetivo general de este trabajo es estudiar el vacío de la
intersección de una gramática de concatenación de rango con un lenguaje
regular finito como marco para la solución del SAT, y construir una
representación de esa intersección sobre la que abordar dicho problema.

Para alcanzar ese objetivo general se plantean los siguientes objetivos
específicos:
\begin{itemize}
  \item Estudiar las gramáticas de concatenación de rango y los
    antecedentes del SAT abordado mediante formalismos de la teoría de
    lenguajes.
  \item Diseñar un algoritmo que decida el vacío de la intersección de una
    gramática de concatenación de rango con un lenguaje regular finito y
    caracterizar su costo.
  \item Modelar el SAT como el vacío de la intersección de una gramática
    de concatenación de rango con un lenguaje regular finito, asociando a
    cada fórmula un autómata booleano y una gramática que impone la
    consistencia entre las ocurrencias de cada variable.
  \item Construir la gramática producto que reconoce la intersección de
    una gramática de concatenación de rango con un lenguaje regular finito,
    demostrar su correctitud y caracterizar su tamaño y el costo de
    construirla.
\end{itemize}

La propuesta de solución parte de un algoritmo de fuerza bruta que decide
el vacío de la intersección de una RCG con un lenguaje regular finito, y
se caracteriza su costo. La consistencia se impone con una RCG de aridad
uno, la mínima posible, gracias a la no linealidad, y se construye la
gramática producto,
cuyo lenguaje es exactamente esa intersección.

Como resultado, se obtiene un marco uniforme para el SAT en torno al
vacío de la intersección de una RCG con un lenguaje regular finito. El
algoritmo de decisión tiene costo exponencial en el peor caso. La
gramática de consistencia tiene aridad uno para cualquier fórmula; a
cambio, se pierde la garantía directa de que el vacío sea polinomial. Se
construye la gramática producto y se demuestra su correctitud; su tamaño
puede ser exponencial, con muchos predicados y cláusulas no productivos.
Queda como problema abierto eliminarlos y analizar el costo de decidir el
vacío sobre la gramática resultante.

Por último, es importante mencionar el uso de la inteligencia
artificial en la elaboración de este trabajo. Como apoyo en la revisión
bibliográfica, se usó para buscar y organizar información en bases de
datos académicas, y así se accedió a bibliografía actualizada y
pertinente para el marco teórico. La opinión de los tutores
presentada en este documento fue generada (por uno de los tutores, con
el consentimiento del otro) usando modelos
grandes de lenguaje, en particular DeepSeek. Para ello se usó el resumen
de la tesis y las características personales del estudiante que resultaron
relevantes para el desarrollo del trabajo. Para alcanzar el tono de
escritura propio de los tutores, se proporcionaron como ejemplos los
mismos elementos (resumen de la tesis y características del estudiante) de
opiniones del tutor de cursos anteriores.

La estructura de este documento es la siguiente. En el
capítulo~\ref{cap:preliminares} se presentan los preliminares de la
teoría de lenguajes, los autómatas finitos, la complejidad computacional
y el SAT. En el capítulo~\ref{cap:rcg-sat} se presentan las gramáticas de
concatenación de rango y los antecedentes del SAT mediante formalismos de
la teoría de lenguajes. En el capítulo~\ref{cap:vacio-interseccion-rcg-regular}
se define el problema del vacío de la intersección de una RCG con un
lenguaje regular finito y se da un algoritmo de decisión. En el
capítulo~\ref{cap:modelacion-sat} se modela el SAT como ese vacío, con una
RCG de consistencia para cualquier fórmula. En el
capítulo~\ref{cap:interseccion} se construye la gramática producto y se
caracterizan su tamaño y el costo de construirla. Por último, en las
conclusiones y recomendaciones se recogen los resultados y las líneas de
investigación que se proponen.

\bibliographystyle{plain}
\bibliography{../bib/referencias}

\end{document}
